\chapter{FUNDAMENTAÇÃO TEÓRICA E TRABALHOS RELACIONADOS}
\label{cap:fundamentacao}

Este capítulo reúne os conceitos necessários para especificar o simulador: a epidemiologia da dengue, a modelagem baseada em agentes e os modelos de estados, a representação espacial, a estocasticidade, a plataforma GAMA e as práticas de documentação, verificação e validação. Ao final, compara estudos relacionados e situa a contribuição deste trabalho.

\section{Dengue, vetor e transmissão}

A dengue é causada pelo vírus da dengue, que possui quatro sorotipos (DENV-1 a DENV-4), e é transmitida principalmente por mosquitos do gênero \textit{Aedes}. A infecção pode ser assintomática ou variar em gravidade, e a circulação de vários sorotipos afeta a resposta imunológica e o risco de formas graves \cite{WHO2023}. Um modelo que trate a doença como uma única dinâmica genérica é uma simplificação aceitável, desde que seja declarada.

No contexto urbano brasileiro, o \textit{Aedes aegypti} é o principal vetor. Somente as fêmeas adultas transmitem o vírus, mas o ciclo de vida inclui os estágios de ovo, larva e pupa, cujo desenvolvimento depende da temperatura e da disponibilidade de água em criadouros. A precipitação e a temperatura influenciam, assim, a população de mosquitos adultos.

A transmissão ocorre em dois sentidos. Um mosquito suscetível adquire o vírus ao picar uma pessoa infectada e só se torna capaz de transmiti-lo após o período de incubação extrínseca. Do mesmo modo, uma pessoa suscetível adquire a infecção pela picada de um mosquito infectado. Como infecção e infectividade não coincidem, os modelos usam um estado intermediário (exposto) para separar os agentes recém-infectados dos que já transmitem. Em modelos espaciais, a ocorrência da transmissão depende ainda do contato entre hospedeiro e vetor, que é determinado pela localização e pela mobilidade dos agentes.

Medidas de controle, como a redução de criadouros e as intervenções sobre o vetor ou sobre a população humana (por exemplo, a vacinação), atuam em componentes diferentes do sistema. Quando representadas em um modelo, cada medida deve ser ligada à regra que corresponde ao seu mecanismo de ação.

\section{Modelagem e simulação computacional}

Um modelo é uma representação simplificada de um sistema, construída para investigar certos aspectos do seu funcionamento. A simulação executa as regras do modelo ao longo do tempo. É útil distinguir o fenômeno real, o modelo conceitual, a implementação e os resultados da simulação, pois divergências entre o esperado e o observado podem vir de qualquer uma dessas camadas.

A Tabela~\ref{tab:abordagens} compara as três abordagens mais comuns. Elas não são excludentes: um modelo baseado em agentes pode adotar estados de modelos compartimentais e produzir indicadores agregados.

\begin{table}[htb]
    \centering
    \caption{Abordagens de simulação de sistemas dinâmicos}
    \label{tab:abordagens}
    \begin{tabular}{p{3.6cm} p{4.6cm} p{5.6cm}}
        \hline
        \textbf{Abordagem} & \textbf{Unidade de representação} & \textbf{Característica principal} \\
        \hline
        Dinâmica de sistemas & Variáveis agregadas (estoques e fluxos) & Relações entre quantidades agregadas ao longo do tempo \\
        Eventos discretos & Eventos que alteram o estado do sistema & Mudanças de estado desencadeadas por eventos \\
        Baseada em agentes & Entidades individuais e ambiente & Heterogeneidade e interações locais, com padrões coletivos emergentes \\
        \hline
    \end{tabular}
    \fonte{Elaborado pelo autor, a partir de \citeonline{Macal2010}.}
\end{table}

A finalidade do modelo determina suas abstrações, seus dados e os critérios de avaliação. Modelos explicativos investigam mecanismos capazes de produzir certos comportamentos; modelos preditivos estimam o comportamento futuro ou desconhecido de um sistema. Reproduzir um padrão observado não prova capacidade preditiva, e um modelo explicativo não precisa fazer previsões para cumprir seu objetivo.

\section{Modelagem baseada em agentes e modelos de estados}

Na modelagem baseada em agentes (ABM), o sistema é representado por entidades com atributos, variáveis de estado e comportamentos próprios, que interagem entre si e com o ambiente. O comportamento agregado emerge dessas interações individuais, sem ser definido por uma regra global \cite{Macal2010}. A abordagem é adequada a sistemas com entidades heterogêneas e interações dependentes do contexto local, como a dengue, em que pessoas e mosquitos têm posições e estados distintos. A emergência de um padrão, porém, não comprova que o modelo reproduz o sistema real.

Em modelos espaciais, a interação costuma depender de regras de vizinhança ou proximidade. Buscar vizinhos exaustivamente em grandes populações exige muitas comparações, de modo que a estratégia de busca afeta o custo computacional.

Os modelos epidemiológicos compartimentais descrevem a evolução da doença por estados. O SIR usa os estados suscetível, infectado e recuperado; o SEIR acrescenta o estado exposto (período entre a infecção e a capacidade de transmitir). Em um ABM, cada agente carrega seu próprio estado e muda de estado por regras, temporizadores ou probabilidades. Para este trabalho, adota-se a estrutura SEIR para humanos e SEI para mosquitos (Figura~\ref{fig:seir}):

\begin{itemize}
    \item \textbf{Suscetível (S):} pode adquirir a infecção;
    \item \textbf{Exposto (E):} infectado, ainda sem capacidade de transmitir;
    \item \textbf{Infectado (I):} capaz de participar da transmissão;
    \item \textbf{Recuperado (R):} concluiu o processo representado pelo modelo (apenas humanos).
\end{itemize}

\begin{figure}[htb]
    \centering
    \caption{Estados epidemiológicos de humanos (SEIR) e mosquitos (SEI)}
    \label{fig:seir}
    \begin{tikzpicture}[node distance=1.6cm, every node/.style={draw, rounded corners, minimum width=1.1cm, minimum height=0.8cm}, >=stealth]
        \node (hs) {S};
        \node[right=of hs] (he) {E};
        \node[right=of he] (hi) {I};
        \node[right=of hi] (hr) {R};
        \draw[->] (hs) -- (he);
        \draw[->] (he) -- (hi);
        \draw[->] (hi) -- (hr);
        \node[draw=none, left=0.3cm of hs, minimum width=0pt] {Humanos};
        \node[below=1.6cm of hs] (ms) {S};
        \node[right=of ms] (me) {E};
        \node[right=of me] (mi) {I};
        \draw[->] (ms) -- (me);
        \draw[->] (me) -- (mi);
        \node[draw=none, left=0.3cm of ms, minimum width=0pt] {Mosquitos};
        \draw[->, dashed] (mi.north) to[bend right=25] node[draw=none, midway, above, font=\scriptsize, minimum height=0pt] {picada} (he.south);
        \draw[->, dashed] (hi.south) to[bend left=25] node[draw=none, midway, below, font=\scriptsize, minimum height=0pt] {picada} (me.north);
    \end{tikzpicture}
    \fonte{Elaborado pelo autor.}
\end{figure}

A implementação de um modelo epidemiológico equivale a uma máquina de estados: é preciso definir os estados válidos, as transições possíveis e as condições de cada transição. Essa especificação facilita a verificação, pois permite comparar o comportamento esperado com o implementado. Atributos como a vacinação podem ser mantidos separados dos estados epidemiológicos, o que explicita como eles influenciam a infecção.

\section{Representação espacial e mobilidade}

Os Sistemas de Informação Geográfica (SIG) organizam dados associados a localizações, em formato vetorial (pontos, linhas e polígonos) ou matricial (superfícies e variáveis distribuídas). Em uma simulação, as geometrias podem delimitar o ambiente, orientar a distribuição dos agentes, restringir deslocamentos e apoiar cálculos de distância e proximidade, além de visualizar o território. Taillandier e Drogoul \cite{taillandier_drogoul_2010} descrevem três usos: ler arquivos espaciais para inicializar ou delimitar o ambiente, usar geometrias como elementos que condicionam operações do modelo e representar objetos geográficos como agentes. Para integrar dados e modelo, é preciso compatibilizar sistemas de coordenadas, unidades, qualidade das geometrias e nível de agregação.

A mobilidade influencia as oportunidades de interação. Escolher um destino e definir a trajetória até ele são problemas distintos. Dados agregados, como os de deslocamento entre regiões, não fornecem as trajetórias individuais, de modo que a mobilidade em um ABM é sempre uma hipótese de modelagem que depende dos objetivos e dos dados disponíveis.

\section{Estocasticidade e experimentação}

Modelos estocásticos incluem processos aleatórios, como a ocorrência de transmissões, o movimento dos agentes e o ciclo de vida dos vetores. Por isso, execuções com os mesmos parâmetros podem produzir resultados diferentes, e a análise precisa considerar essa variabilidade.

A semente do gerador pseudoaleatório controla a sequência de valores sorteados. Com a mesma semente, configuração, condições iniciais e versão da implementação, é possível, em geral, reproduzir uma execução. Uma única execução pode não representar a variabilidade do modelo; por isso, usam-se múltiplas replicações e medidas descritivas, como média e dispersão. O número de replicações depende dos objetivos e do custo computacional, e diferenças entre cenários só são relevantes se forem maiores que a variabilidade entre replicações.

A análise de sensibilidade investiga como mudanças nos parâmetros afetam as saídas, por variação controlada de parâmetros ou por métodos globais, como os índices de Sobol, que estimam a contribuição de cada parâmetro para a variância das saídas \cite{sobol_2001}. Em modelos espaciais, a incerteza também pode variar no território e no tempo \cite{gamal_etal_2025}.

\section{Plataforma GAMA e linguagem GAML}

O GAMA é uma plataforma para construir e executar modelos espaciais baseados em agentes, com recursos para agentes, ambientes geográficos e experimentos \cite{Taillandier2019}. Seus modelos são escritos em GAML, linguagem que define variáveis globais, espécies de agentes, inicialização, comportamentos e experimentos. A plataforma oferece operações espaciais, como cálculo de distâncias, identificação de vizinhos, manipulação de geometrias e definição de trajetórias.

A plataforma facilita a integração entre agentes e dados geográficos, mas não garante a qualidade do modelo. A descrição de um simulador deve indicar quais recursos da plataforma foram efetivamente usados.

\section{Documentação, verificação e validação}

O protocolo ODD (\textit{Overview, Design concepts and Details}) padroniza a documentação de ABMs, organizando propósito, entidades, escalas, processos, inicialização, entradas e submodelos \cite{grimm_etal_2020}. Documentar torna as decisões explícitas e facilita a reprodução, mas não prova que o código executa as regras especificadas.

Três conceitos devem ser distinguidos:

\begin{itemize}
    \item \textbf{Verificação:} analisa se a implementação corresponde à especificação, por meio de testes de transições, inicialização, regras de interação e indicadores. É necessária porque uma simulação pode parecer plausível mesmo com regras implementadas de modo diferente do especificado.
    \item \textbf{Validação:} avalia se o modelo é adequado ao seu propósito e representa aspectos relevantes do sistema real, o que pode incluir comparação com dados observados. Um código correto pode ainda usar hipóteses insuficientes, e a semelhança visual com uma série observada não comprova os mecanismos internos.
    \item \textbf{Calibração:} ajusta parâmetros com base em dados observados.
\end{itemize}

Dois requisitos sustentam essas atividades. A rastreabilidade liga os elementos conceituais (entidades, estados, regras) aos componentes do código. A reprodutibilidade exige registrar parâmetros, sementes, versões de software e procedimentos experimentais, o que é essencial em modelos estocásticos.

\section{Trabalhos relacionados}

Os estudos a seguir servem de referência para decisões de modelagem, representação espacial e experimentação. Como cada um tem área, hipóteses e dados próprios, seus parâmetros e resultados não se transferem automaticamente a este trabalho. A Tabela~\ref{tab:relacionados} resume as contribuições.

\begin{table}[htb]
    \centering
    \caption{Comparação dos trabalhos relacionados}
    \label{tab:relacionados}
    \begin{tabular}{p{3.0cm} p{2.7cm} p{2.8cm} p{5.1cm}}
        \hline
        \textbf{Estudo} & \textbf{Doença e área} & \textbf{Plataforma} & \textbf{Contribuição para este trabalho} \\
        \hline
        \citeonline{garcias_2024} & Dengue, Portugal Continental & GAMA & Documentação do modelo; relação entre espaço, regras e experimentação \\
        \citeonline{oliveira_etal_2025} & Dengue, Alto Santo (CE) & GAMA & Organização de entidades (humanos, mosquitos, ovos, criadouros), parâmetros e cenários \\
        \citeonline{gharakhanlou_2019, gharakhanlou_2020} & Malária & Não aplicável & Integração entre entidades, espaço, ambiente e intervenções \\
        \citeonline{gamal_etal_2025} & Não específica de doença & Não aplicável & Análise espaço-temporal da incerteza em ABM \\
        \hline
    \end{tabular}
    \fonte{Elaborado pelo autor.}
\end{table}

Garcias \cite{garcias_2024} modela a dengue em Portugal Continental no GAMA, com agentes humanos, mosquitos e componentes territoriais, e discute parâmetros, mobilidade e condições ambientais. Oliveira et al. \cite{oliveira_etal_2025} simulam a dengue em Alto Santo (CE), também no GAMA, com humanos, mosquitos, ovos e criadouros, informações espaciais e experimentos com diferentes parâmetros. Em ambos os casos, as diferenças de espécie vetorial, área de estudo e condições epidemiológicas impedem reutilizar os parâmetros em São Paulo.

Os modelos de malária de Gharakhanlou et al. \cite{gharakhanlou_2019, gharakhanlou_2020} trazem aspectos metodológicos sobre representação espacial, mas suas regras epidemiológicas não se aplicam à dengue. Gamal et al. \cite{gamal_etal_2025} propõem uma forma de analisar a propagação espaço-temporal da incerteza em ABMs, mostrando que a variabilidade dos resultados difere entre unidades espaciais e ao longo do tempo. É uma referência conceitual para avaliação, e não um conjunto de regras para a dengue.

Em conjunto, os trabalhos mostram como documentar modelos, representar entidades, integrar dados geográficos e conduzir experimentos. Os estudos de dengue descritos aqui, porém, tratam de outros territórios. Este trabalho aplica a abordagem ao município de São Paulo e acrescenta a verificação sistemática do simulador e a análise da variabilidade estocástica, com a especificação e a avaliação próprias descritas nos capítulos seguintes.